\section{Riemann--Stieltjes integrals}
\label{sec:rs-integrals}

\textbf{Read this before calculating.}
An ordinary integral weights an interval by its length. A Stieltjes
integral weights it by the increase of an increasing function $F$.
Here ``increasing'' permits equality: $s<t$ implies $F(s)\le F(t)$.
Ross uses a modified definition that handles jumps explicitly.
This matters: a formula from another book's traditional Stieltjes
theory may not apply to this section.

Write $F(t-)$ and $F(t+)$ for the one-sided limits. On the particular
interval $[a,b]$ being integrated over, the endpoint conventions are
$F(a-)=F(a)$ and $F(b+)=F(b)$. Thus the mass at $t$ is
$\Delta F(t)=F(t+)-F(t-)$. For a partition
$P=\{a=t_0<\cdots<t_n=b\}$ set
\[
 w_k=F(t_k-)-F(t_{k-1}+),\qquad
 J_F(f,P)=\sum_{k=0}^n f(t_k)\Delta F(t_k).
\]
The upper sum is $J_F(f,P)+\sum_{k=1}^n (\sup_{(t_{k-1},t_k)}f)w_k$;
the lower sum uses infima instead. The numbers $w_k$ and
$\Delta F(t_k)$ are nonnegative, and together they add to
$W=F(b)-F(a)$. In particular, if $\abs{h}\le B$, then every upper
and lower sum for $h$ lies between $-BW$ and $BW$.

The essential test is Theorem~35.6: a bounded function is
$F$-integrable precisely when, for every $\varepsilon>0$, some
partition satisfies $U_F(f,P)-L_F(f,P)<\varepsilon$. The jump terms
cancel in this difference. We will often use the equivalent expression
\[
 U_F(f,P)-L_F(f,P)
 =\sum_{k=1}^n\operatorname{osc}(f,(t_{k-1},t_k))w_k,
\]
where $\operatorname{osc}(f,I)=\sup_I f-\inf_I f$ is the oscillation
of $f$ on $I$.

\exercise{35.1}{Why an increasing function has one-sided limits}
\goal Identify both one-sided limits of $F$ by a supremum or an
infimum, including the indicated endpoints.
\idea Increasing values cannot move back and forth. Their supremum
on the left, or infimum on the right, is therefore the only possible
limiting value. We verify this using the definition of a limit.

\step{(a) Approach $t\in(a,b]$ from the left.}
The set $A=\{F(x):a<x<t\}$ is nonempty and bounded above by $F(t)$.
Consequently $L=\sup A$ is a finite real number. Given
$\varepsilon>0$, the definition of supremum supplies $y\in(a,t)$
such that $L-\varepsilon<F(y)\le L$. For every $x$ with $y<x<t$,
monotonicity gives
\[
 L-\varepsilon<F(y)\le F(x)\le L.
\]
Take $\delta=t-y>0$. If $x\in[a,b]$ and
$0<t-x<\delta$, then $y<x<t$, so $\abs{F(x)-L}<\varepsilon$.
This is exactly $\lim_{x\to t-}F(x)=L$. The proof works when
$t=b$, since it only uses points to the left of $t$.

\step{(b) Approach $t\in[a,b)$ from the right.}
The nonempty set $B=\{F(x):t<x<b\}$ is bounded below by $F(t)$.
Put $R=\inf B$. Given $\varepsilon>0$, choose $y\in(t,b)$
with $R\le F(y)<R+\varepsilon$. Whenever $t<x<y$,
\[
 R\le F(x)\le F(y)<R+\varepsilon.
\]
Now $\delta=y-t$ proves $\lim_{x\to t+}F(x)=R$.
At $t=a$ the same argument is valid using only points to the right.
\answer{The left limit is $\sup\{F(x):a<x<t\}$ and the right
limit is $\inf\{F(x):t<x<b\}$. Both are finite.}

\exercise{35.2}{Separate a continuous increase from a jump}
\goal Evaluate $\int_0^3x^2\,dF(x)$ for the piecewise function in
Example~4.
\step{Record all three pieces, especially the jump.}
The function is $F(t)=0$ for $t<0$, $F(t)=t^2$ for $0\le t<2$,
and $F(t)=t+5$ for $t\ge2$. At $2$ its left limit is $4$ and its
right limit is $7$, so the jump has size $3$.
Write $F=F_1+3J_2$, where on $[0,3]$
\[
 F_1(t)=\begin{cases}t^2,&0\le t\le2,\\t+2,&2\le t\le3.\end{cases}
\]
This $F_1$ is continuous. Its derivative is $2t$ on the first
piece and $1$ on the second. We apply Theorem~35.13 separately on
the two pieces; no derivative at the joining point is needed.

\step{Calculate each contribution.}
Linearity in the integrator, interval additivity, and the jump rule
$\int f\,dJ_2=f(2)$ give
\begin{align*}
 \int_0^3x^2\,dF(x)
 &=\int_0^2x^2(2x)\,dx+\int_2^3x^2\,dx+3(2^2)\\
 &=\left[\frac{x^4}{2}\right]_0^2
   +\left[\frac{x^3}{3}\right]_2^3+12\\
 &=8+\frac{27-8}{3}+12=\frac{79}{3}.
\end{align*}
The extra $12$ is essential: differentiating the smooth pieces alone
would lose the mass concentrated at $2$.
\answer{$\displaystyle\int_0^3x^2\,dF(x)=79/3$.}

\exercise{35.3}{Integrating against the integer step function}
\goal Evaluate the three requested integrals against $F(t)=\lfloor t\rfloor$.
\idea This function has jumps of size $1$ at the integers. On a
specified integration interval we include only the jumps belonging
to the restricted function, with Ross's endpoint convention.

\step{(a) Integrate $x$ on $[0,6]$.}
On this interval $F(0-)$ means $F(0)=0$, rather than the limit of
the original function from negative numbers. Also $F(0+)=0$.
There is therefore no mass at $0$. There is a mass of $1$ at
each of $1,2,3,4,5,6$. The last one is included because
$F(6+)=F(6)=6$ and $F(6-)=5$. Example~1 gives
\[
 \int_0^6x\,dF(x)=\sum_{j=1}^6j=\frac{6\cdot7}{2}=21.
\]

\step{(b) Integrate $x^2$ on $[0,3]$.}
The contributing jumps are at $1,2,3$, with the same endpoint
reasoning. Thus
\[
 \int_0^3x^2\,dF(x)=1^2+2^2+3^2=1+4+9=14.
\]

\step{(c) Integrate $x^2$ on $[1/4,\pi/4]$.}
Both endpoints lie strictly between $0$ and $1$. The integrator is
identically $0$ on the whole interval; every interval increment and
every endpoint jump is zero. Its upper and lower sums are both zero.
\answer{The three values are $21$, $14$, and $0$, respectively.}

\exercise{35.4}{A smooth trigonometric integrator}
\goal Integrate $x$ against $F(t)=\sin t$ on the two specified intervals.
\step{Check the conversion to an ordinary integral.}
On $[-\pi/2,\pi/2]$, $F'(t)=\cos t\ge0$, so $F$ is increasing.
Both $F'$ and $f(t)=t$ are continuous. Theorem~35.13 therefore
permits the replacement $dF(x)=\cos x\,dx$ on each interval.

\step{(a) The interval $[0,\pi/2]$.}
Integration by parts with $u=x$ and $dv=\cos x\,dx$ gives
\[
 \int x\cos x\,dx=x\sin x-\int\sin x\,dx
                  =x\sin x+\cos x.
\]
Consequently
\[
 \int_0^{\pi/2}x\,dF(x)
 =\left[x\sin x+\cos x\right]_0^{\pi/2}
 =\frac\pi2-1.
\]

\step{(b) The interval $[-\pi/2,\pi/2]$.}
The function $x\cos x$ is odd: replacing $x$ by $-x$ changes its
sign. Its integrals over the two halves cancel, so the answer is
$0$. Alternatively the displayed antiderivative has value $\pi/2$
at both endpoints, giving difference $0$.
\answer{The values are $\pi/2-1$ and $0$.}

\exercise{35.5}{The rational indicator sees continuous and atomic mass differently}
\goal For $f=\mathbf{1}_{\Q}$, contrast continuous integrators
with the purely discrete integrators in Examples~1 and~3.

\step{(a) A continuous nonconstant integrator.}
Every nonempty open interval contains both a rational number and
an irrational number. Hence on every partition interval
$(t_{k-1},t_k)$, the supremum of $f$ is $1$ and the infimum is $0$.
Continuity of $F$ makes all jump terms zero and gives
$w_k=F(t_k)-F(t_{k-1})$. Therefore, for every partition,
\[
 L_F(f,P)=0,\qquad
 U_F(f,P)=\sum_{k=1}^n[F(t_k)-F(t_{k-1})]=F(b)-F(a).
\]
Taking the supremum of the lower sums and infimum of the upper
sums changes neither number. Since $F(b)>F(a)$, the lower and
upper integrals are different. Thus $f$ is not $F$-integrable.

\step{(b) A finite or summable collection of point masses.}
For $F=\sum_jc_jJ_{u_j}$ as in those examples, every bounded
function is $F$-integrable, including this $f$. Its integral is
\[
 \int_a^bf\,dF=\sum_jc_jf(u_j)
              =\sum_{\{j:u_j\in\Q\}}c_j.
\]
The formula includes the finite case. In the infinite case it
converges absolutely, because $0\le c_jf(u_j)\le c_j$ and
$\sum_jc_j<\infty$. These integrators assign their mass to actual
points; the value of $f$ at each mass point is well defined.
\answer{For continuous $F$ of positive total increase, $f$ is not
integrable. For the discrete $F$ of Examples~1 and~3, it is
integrable, with integral equal to the total mass at rational points.}

\exercise{35.6}{Uniform limits preserve Stieltjes integrability}
\goal If each $f_n$ is $F$-integrable and $f_n\to f$ uniformly on
$[a,b]$, prove that $f$ is $F$-integrable and that the integrals converge.
\idea Uniform convergence controls every point by the same error.
Multiplying that error by the total mass $W=F(b)-F(a)$ controls
every upper sum, lower sum, and ultimately the integral.

\step{First establish boundedness.}
Take one index $N$ for which $\abs{f(x)-f_N(x)}<1$ for every
$x\in[a,b]$. Such an index exists by uniform convergence. Since
$f_N$ is integrable, it is bounded, say $\abs{f_N(x)}\le B$.
Then $\abs{f(x)}\le B+1$ everywhere. This verifies the boundedness
required before applying the Darboux criterion.

\step{Make the upper--lower gap small.}
Assume $W>0$; if $W=0$, every sum and integral is zero and the
conclusions are immediate. Fix $\varepsilon>0$ and select $N$ so
that $\abs{f-f_N}\le\eta$ everywhere, where
$\eta=\varepsilon/(4W)$. On any nonempty interval $I$,
\[
 \sup_I f\le\sup_I f_N+\eta,\qquad
 \inf_I f\ge\inf_I f_N-\eta.
\]
Subtracting gives
$\operatorname{osc}(f,I)\le\operatorname{osc}(f_N,I)+2\eta$.
Because $f_N$ is integrable, choose $P$ with
$U_F(f_N,P)-L_F(f_N,P)<\varepsilon/2$. Since $\sum_kw_k\le W$,
\begin{align*}
 U_F(f,P)-L_F(f,P)
 &\le U_F(f_N,P)-L_F(f_N,P)+2\eta\sum_kw_k\\
 &<\frac\varepsilon2+2\eta W=\varepsilon.
\end{align*}
Theorem~35.6 proves that $f$ is integrable.

\step{Now estimate the integrals.}
Put $\eta_n=\sup_{x\in[a,b]}\abs{f_n(x)-f(x)}$; uniform
convergence says $\eta_n\to0$. The integrable difference
$f_n-f$ satisfies $-\eta_n\le f_n-f\le\eta_n$. Order preservation
and integration of constants give
\[
 -\eta_nW\le\int_a^b(f_n-f)\,dF\le\eta_nW.
\]
Thus $\abs{\int_a^bf_n\,dF-\int_a^bf\,dF}\le\eta_nW\to0$.
\answer{Uniform convergence permits passage to the limit under this
integral, with error at most the uniform function error times the total mass.}

\exercise{35.7}{Products, squares, maxima, and minima}
\goal Show that all four indicated constructions preserve $F$-integrability.

\step{(a) Squaring an integrable function.}
Choose $B\ge1$ such that $\abs{f(x)}\le B$ on $[a,b]$.
For any two points $x,y$ in the same interval $I$,
\[
 \abs{f(x)^2-f(y)^2}
 =\abs{f(x)-f(y)}\abs{f(x)+f(y)}
 \le2B\abs{f(x)-f(y)}.
\]
The supremum of pairwise differences of a bounded real function
is its oscillation. It follows that
$\operatorname{osc}(f^2,I)\le2B\operatorname{osc}(f,I)$.
Given $\varepsilon>0$, take a partition making the gap for $f$
less than $\varepsilon/(2B)$. The gap formula then makes the gap
for $f^2$ less than $\varepsilon$. Since $f^2$ is bounded, the
criterion proves its integrability.

\step{(b) Multiplying two integrable functions.}
Linearity makes $f+g$ and $f-g$ integrable. Part (a) makes their
squares integrable. The identity
\[
 fg=\frac{(f+g)^2-(f-g)^2}{4}
\]
then expresses $fg$ as a linear combination of integrable functions.
To check the identity, expand both squares: their $f^2$ and $g^2$
terms cancel, leaving $4fg$ in the numerator.

\step{(c) Taking the pointwise maximum and minimum.}
Theorem~35.10 makes $\abs{f-g}$ integrable. At each point,
\[
 \max(f,g)=\frac{f+g+\abs{f-g}}2,\qquad
 \min(f,g)=\frac{f+g-\abs{f-g}}2.
\]
For example, if $f(x)\ge g(x)$ then $\abs{f(x)-g(x)}=f(x)-g(x)$,
so the first expression becomes $f(x)$ and the second becomes
$g(x)$. The other ordering gives the reversed roles. The formulas
and linearity prove both integrability assertions.
\answer{$f^2$, $fg$, $\max(f,g)$, and $\min(f,g)$ are all $F$-integrable.}

\exercise{35.8}{A continuous density defines the weight}
\goal For continuous $g\ge0$ and $F(t)=\int_a^tg(x)\,dx$, show
that $\int_a^bf\,dF=\int_a^bf(x)g(x)\,dx$ for continuous $f$.

\step{Check that $F$ is an admissible increasing integrator.}
If $a\le s<t\le b$, interval additivity gives
\[
 F(t)-F(s)=\int_s^tg(x)\,dx\ge0,
\]
because $g$ is nonnegative. Thus $F$ is increasing. The Fundamental
Theorem of Calculus, applied to continuous $g$, gives $F'=g$ on
the interior, with the corresponding one-sided derivatives at the
endpoints. In particular $F'$ is continuous on the closed interval.

\step{Apply the smooth-integrator formula.}
All the hypotheses of Theorem~35.13 are now verified, so
\[
 \int_a^bf(x)\,dF(x)
 =\int_a^bf(x)F'(x)\,dx
 =\int_a^bf(x)g(x)\,dx.
\]
If $g$ is identically zero, $F$ is constant. Although the section
usually excludes a constant integrator to avoid trivial cases, its
sums are still defined and all equal zero; the formula is then also
true. If $g$ is nonzero somewhere, continuity and nonnegativity
give positive integral, so $F(b)>F(a)$ as in the standing convention.
\answer{The density $g$ converts a Stieltjes integral into the usual
weighted Riemann integral.}

\exercise{35.9}{A weighted mean value is attained}
\goal Prove the mean-value assertion for continuous $f$, then
recover the weighted integral theorem in Exercise~33.14.

\step{(a) Trap the average between a minimum and a maximum.}
The Extreme Value Theorem gives real numbers
$m=\min_{[a,b]}f$ and $M=\max_{[a,b]}f$. Put $W=F(b)-F(a)>0$.
The continuous function $f$ is $F$-integrable by Theorem~35.7.
Integrating $m\le f\le M$ with order preservation gives
\[
 mW\le\int_a^bf\,dF\le MW.
\]
Divide by the positive number $W$ and call the quotient $A$.
We have $m\le A\le M$. By the Intermediate Value Theorem,
$f$ takes every value between its attained minimum and maximum.
Hence some $x\in[a,b]$ satisfies $f(x)=A$. Multiplication by $W$
proves
\[
 \int_a^bf\,dF=f(x)[F(b)-F(a)].
\]
The conclusion naturally allows endpoints: all the mass could be
concentrated at $a$, for example.

\step{(b) Recover the continuous-density theorem.}
For continuous $g\ge0$ set $F(t)=\int_a^tg(s)\,ds$, as in
Exercise~35.8. If $g\equiv0$, both integrals in Exercise~33.14(a)
are zero and any $x\in(a,b)$ works. Otherwise $W=\int_a^bg>0$.
Part (a) and Exercise~35.8 give
\[
 \int_a^bf(t)g(t)\,dt=f(x)\int_a^bg(t)\,dt
\]
initially for some $x\in[a,b]$.

\step{Explain why the point can be chosen in the interior here.}
Exercise~33.14(a) asks for $x\in(a,b)$, so this small extra point
deserves a proof. Let
$A=(\int_a^bfg)/(\int_a^bg)$. Suppose no interior point satisfies
$f(x)=A$. By continuity and the Intermediate Value Theorem,
$f-A$ must then have the same strict sign on all of $(a,b)$.
Also some $t_0\in(a,b)$ has $g(t_0)>0$: a positive endpoint
value of $g$ would imply positive nearby interior values as well.
If $f-A>0$, continuity gives a short interval around $t_0$ on which
both $f-A$ and $g$ have positive lower bounds. Thus
$(f-A)g\ge0$ everywhere and has strictly positive integral, contrary to
\[
 \int_a^b(f-A)g=\int_a^bfg-A\int_a^bg=0.
\]
If $f-A<0$, apply the same argument to $(A-f)g$ and obtain the
same contradiction. Therefore an interior point does attain $A$.
Taking $g\equiv1$ recovers the ordinary integral mean-value
theorem, as requested in Exercise~33.14(b).
\answer{Part (a) is the Stieltjes mean-value formula. With a
continuous nonnegative density it gives Exercise~33.14(a), including
its interior-point conclusion.}

\exercise{35.10}{A polynomial integrated against itself}
\goal For a positive odd integer $n$ and $F(x)=f(x)=x^n$, evaluate
$\int_a^bf\,dF$.
\step{Explain the parity hypothesis.}
The map $x\mapsto x^n$ is increasing on all of $\R$ when $n$ is
odd. Its derivative is $nx^{n-1}\ge0$ since $n-1$ is even.
Both $f$ and $F'$ are continuous, so Theorem~35.13 applies on
any closed interval $[a,b]$.
\step{Convert and integrate.}
We have
\begin{align*}
 \int_a^bx^n\,d(x^n)
 &=\int_a^bx^n\,nx^{n-1}\,dx
   =n\int_a^bx^{2n-1}\,dx\\
 &=n\left[\frac{x^{2n}}{2n}\right]_a^b
   =\frac12(b^{2n}-a^{2n}).
\end{align*}
The result can be negative if the interval lies on the negative
half-line. That is consistent: the weight is nonnegative, but
$f(x)=x^n$ is negative there.
\answer{$\displaystyle\int_a^bf\,dF=\tfrac12[b^{2n}-a^{2n}]$.}

\exercise{35.11}{Change variables without differentiating the substitution}
\goal For a continuous strictly increasing map
$\phi:[c,d]\to[a,b]$ with the stated endpoint values, prove
integrability and equality of the two Stieltjes integrals for
$g=f\circ\phi$ and $G=F\circ\phi$.
\idea A partition and its image under $\phi$ describe exactly the
same intervals and mass points. Match their sums term by term.
There is no assumption that $\phi$ is differentiable.

\step{The substitution maps intervals onto intervals.}
Continuity, strict increase, and the endpoint values imply that
$\phi$ is a bijection from $[c,d]$ onto $[a,b]$. Its inverse is
continuous: for an interior $u$, points just below and above $u$
have image values strictly below and above $\phi(u)$, which trap
the inverse of every sufficiently nearby image point. The same
one-sided trapping works at the endpoints. In particular,
\[
 \phi((u,v))=(\phi(u),\phi(v))\qquad(c\le u<v\le d).
\]
The composition $G$ is increasing, and $g$ is bounded because
$f$ is bounded on the image $[a,b]$.

\step{Match one-sided limits and jumps.}
For $u\in(c,d)$, as $v\to u-$ the numbers $\phi(v)$ approach
$\phi(u)$ from below. They cover a whole one-sided interval near
$\phi(u)$. Therefore
\[
 G(u-)=F(\phi(u)-),\qquad G(u+)=F(\phi(u)+).
\]
At $u=c,d$ the endpoint conventions give these same identities
for the relevant decreed limits. Subtracting proves
$\Delta G(u)=\Delta F(\phi(u))$.

\step{Match the upper and lower sums.}
Take any $Q=\{c=u_0<\cdots<u_n=d\}$ and its image partition
$P=\{a=\phi(u_0)<\cdots<\phi(u_n)=b\}$. Bijection on the open
subintervals gives
\[
 \sup_{(u_{k-1},u_k)}g
 =\sup_{(\phi(u_{k-1}),\phi(u_k))}f,
\]
and the analogous identity for infima. The one-sided-limit
identities give matching interval weights, and $g(u_k)=f(\phi(u_k))$
gives matching jump terms. Thus
\[
 U_G(g,Q)=U_F(f,P),\qquad L_G(g,Q)=L_F(f,P).
\]
Every partition of $[a,b]$ arises this way, by applying $\phi^{-1}$
to its points. Consequently the infima of the upper sums agree,
and the suprema of the lower sums agree. Since $f$ is
$F$-integrable, those two common numbers coincide.
\answer{$g$ is $G$-integrable, and
$\displaystyle\int_c^dg\,dG=\int_a^bf\,dF$.}

\exercise{35.12}{A strictly increasing function with dense jumps}
\goal Analyze the sum of positive jumps placed at every rational
point of $[a,b]$.

\step{First address the wording in the supplied edition.}
The printed exercise says that the $c_j$ are ``positive integers''
and that $\sum_jc_j<\infty$. These conditions are incompatible:
each partial sum is at least its number of terms, so the series
would diverge to $+\infty$. Thus there is no example satisfying
the literal hypotheses. The intended, nonvacuous version uses
\emph{positive real numbers} with finite sum, just as in Example~3.
We solve that version explicitly; for instance $c_j=2^{-j}$ is an
allowed choice. Let $(u_j)$ enumerate the rationals in $[a,b]$
without repetition, and retain Ross's special convention for $J_a$.

\step{(a) Show that the sum exists and increases strictly.}
Since $0\le J_{u_j}(t)\le1$,
\[
 0\le\sum_{j=1}^Nc_jJ_{u_j}(t)\le\sum_{j=1}^{\infty}c_j<\infty.
\]
So the nondecreasing partial sums have a finite limit $F(t)$ at
every point. Also their tails are at most $\sum_{j>N}c_j$,
independently of $t$, so the convergence is uniform.
For $a\le s<t\le b$, density of $\Q$ supplies a rational
$u_m\in(s,t)$. The function $J_{u_m}$ is $0$ at $s$ and $1$ at
$t$. Every other summand is increasing, hence has nonnegative
increment. Subtracting the convergent series gives
\[
 F(t)-F(s)
 =\sum_{j=1}^{\infty}c_j[J_{u_j}(t)-J_{u_j}(s)]
 \ge c_m>0.
\]
This proves strict increase, not merely nondecrease.

\step{(b) Prove continuity at every irrational point.}
Let $t\in[a,b]\setminus\Q$ and $\varepsilon>0$. Choose $N$
with $\sum_{j>N}c_j<\varepsilon$. Because $t$ differs from each
of the finitely many $u_1,\ldots,u_N$, there is $\delta>0$ such
that $\abs{x-t}<\delta$ crosses none of these points. Each of
the first $N$ step functions has the same value at $x$ and $t$.
For the remaining terms, $\abs{J_{u_j}(x)-J_{u_j}(t)}\le1$.
Thus for $x\in[a,b]$ with $\abs{x-t}<\delta$,
\[
 \abs{F(x)-F(t)}
 \le\sum_{j>N}c_j\abs{J_{u_j}(x)-J_{u_j}(t)}
 \le\sum_{j>N}c_j<\varepsilon.
\]
This proves relative continuity at $t$, including an irrational
endpoint if one occurs.

\step{Show discontinuity at every rational point.}
If $u_m\in(a,b]$, then $J_{u_m}(u_m)=1$, whereas
$J_{u_m}(x)=0$ for every $x<u_m$. Since every summand is
increasing, $F(u_m)-F(x)\ge c_m$ for $x<u_m$ in $[a,b]$.
Points of the domain can approach $u_m$ from the left while this
gap stays at least $c_m>0$, so $F$ is discontinuous there.
If the rational point is $u_m=a$, Ross defines $J_a(a)=0$ and
$J_a(x)=1$ for $x>a$. Thus $F(x)-F(a)\ge c_m$ for all $x>a$,
which proves right discontinuity at $a$.
In fact Example~3 shows that the jump at $u_m$ is exactly $c_m$.
\answer{Under the intended summable positive-real-weight
hypothesis, $F$ is strictly increasing, continuous exactly at
irrational points of $[a,b]$, and discontinuous at every rational
point. Under the printed positive-integer hypothesis, no such
weight sequence exists.}
